\begin{helper}
Fix a finite terminal coordinate set \(U\), a number \(p\) with
\(0\le p\le 1/2\), and a fixed blue set \(B\subseteq U\) with
\(|B|=u_B\).  Let \(\mathcal B\) be a finite branch-label set with traces
\(\operatorname{blueTrace}(\beta)\), let \(\mathcal T\) be a finite set of
transcripts \(\tau=(P_\tau,A_\tau,Z_\tau)\), let \(J\) be a fixed red label,
and let \(R_{\beta,J}\) be the red support attached to \((\beta,J)\).  Suppose
there is a realized-cell predicate \(\operatorname{Realized}(\beta,\tau)\).

Assume the following reduced hypotheses.  The set \(B\) is blue.  Branch labels
are functional from their complete blue terminal traces: if two labels have
the same truth value of \(\operatorname{blueTrace}\) on every blue coordinate
of \(U\), then the labels are equal.  For every realized \((\beta,\tau)\) with
\(\tau\in\mathcal T\), the realized-cell geometry gives
\[
|R_{\beta,J}|=u_R,\quad R_{\beta,J}\subseteq U,\quad
R_{\beta,J}\text{ red},\quad Z_\tau\text{ red},\quad
P_\tau,A_\tau\subseteq U,\quad P_\tau\cap A_\tau=\emptyset,
\]
\[
B\cup R_{\beta,J}\subseteq P_\tau,\qquad Z_\tau\subseteq A_\tau,
\]
and for every blue \(e\in U\),
\[
e\in\operatorname{blueTrace}(\beta)\Longleftrightarrow e\in P_\tau,
\qquad
e\notin\operatorname{blueTrace}(\beta)\Longleftrightarrow e\in A_\tau .
\]

Let
\[
L=\{(\beta,\tau):\tau\in\mathcal T,
      \operatorname{Realized}(\beta,\tau)\}.
\]
For \(\ell=(\beta,\tau)\in L\), let \(C_\ell\) be the set of
\(A\subseteq U\) such that
\[
P_\tau\subseteq A,\qquad
A\cap(A_\tau\setminus Z_\tau)=\emptyset,\qquad
\forall\text{ blue }e\in U\setminus Z_\tau,
  e\in\operatorname{blueTrace}(\beta)\Longleftrightarrow e\in A .
\]
Then the complete-assignment product mass of the released union is bounded by
the fixed support factor:
\[
\sum_{A\in\bigcup_{\ell\in L}C_\ell}p^{|A|}(1-p)^{|U|-|A|}
\le p^{u_R}p^{u_B}.
\]
\end{helper}
\begin{proof}
Write \(q=1-p\).  The complete-assignment weight of
\(A\subseteq U\) is \(p^{|A|}q^{|U|-|A|}\).  Since \(0\le p\le 1/2\), both
\(p\) and \(q\) are nonnegative; therefore every product weight below is
nonnegative, and enlarging a family of assignments can only increase the
corresponding sum.  Split the terminal coordinates into the disjoint blue and
red parts
\[
U_B=\{e\in U:e\text{ is blue}\},\qquad
U_R=\{e\in U:e\text{ is red}\}.
\]
For an assignment \(A\), write \(A_B=A\cap U_B\) and \(A_R=A\cap U_R\).  Since
the two coordinate classes are disjoint and cover \(U\), the product weight
factors as
\[
p^{|A|}q^{|U|-|A|}
=\bigl(p^{|A_B|}q^{|U_B|-|A_B|}\bigr)
 \bigl(p^{|A_R|}q^{|U_R|-|A_R|}\bigr).
\]

Consider \(A\in C_\ell\) for a realized leaf \(\ell=(\beta,\tau)\).  The
condition \(P_\tau\subseteq A\), together with
\[
B\cup R_{\beta,J}\subseteq P_\tau,
\]
implies \(B\subseteq A_B\) and \(R_{\beta,J}\subseteq A_R\).  The support
sets are disjoint because \(B\) is blue and \(R_{\beta,J}\) is red, with
\(|B|=u_B\) and \(|R_{\beta,J}|=u_R\).

The released set \(Z_\tau\) is red, so every blue coordinate of \(U\) lies
outside \(Z_\tau\).  Hence the released-cylinder trace condition says, for
every blue \(e\in U\),
\[
e\in A_B\Longleftrightarrow e\in\operatorname{blueTrace}(\beta).
\]
Thus \(A_B\) is the complete blue trace of \(\beta\) on \(U\).  If the same
assignment \(A\) is represented by another realized leaf
\((\beta',\tau')\), the same argument gives equality of the two complete blue
traces on every blue terminal coordinate.  Branch functionality then gives
\(\beta=\beta'\).  Applying
\(\noderef{FixedSetHistoryCellRedSupportDeterminedByBlueTrace}\) with this
blue-trace equality shows
\[
R_{\beta,J}=R_{\beta',J}.
\]
More generally, if \(A\in C_{\beta,\tau}\) and \(A'\in C_{\beta',\tau'}\) have
the same blue part \(A_B=A'_B=S\), then both
\(\operatorname{blueTrace}(\beta)\) and
\(\operatorname{blueTrace}(\beta')\) agree with \(S\) on every blue terminal
coordinate.  The same branch-functionality and noderef argument gives
\(\beta=\beta'\) and \(R_{\beta,J}=R_{\beta',J}\).  Therefore a fixed blue
trace \(S\subseteq U_B\) determines at most one branch label and one red
support \(R_S\) relevant to the released union.

Now decompose the union mass by complete blue trace.  For a fixed blue trace
\(S\subseteq U_B\), all assignments in the corresponding trace class have
blue part \(A_B=S\).  If this trace class contributes to the released union,
then \(B\subseteq S\); otherwise its contribution is zero.  Its blue factor is
\[
p^{|S|}q^{|U_B|-|S|}.
\]
For such a contributing trace \(S\), the previous paragraph gives a
trace-determined red support \(R_S\subseteq U_R\) of size \(u_R\), and every
released assignment in the trace class has \(R_S\subseteq A_R\).  Hence the
red part of the released class is contained in the larger family of all
\(T\subseteq U_R\) with \(R_S\subseteq T\).  Summing over this larger family,
the free red coordinates in \(U_R\setminus R_S\) contribute independent
factors \(p+q=1\), so
\[
\sum_{\substack{T\subseteq U_R\\ R_S\subseteq T}}
  p^{|T|}q^{|U_R|-|T|}
=p^{|R_S|}=p^{u_R}.
\]
Therefore the total released mass in the trace class \(S\) is at most
\[
p^{|S|}q^{|U_B|-|S|}\,p^{u_R}.
\]

Summing this bound over all contributing blue traces and then enlarging to
all traces \(S\subseteq U_B\) with \(B\subseteq S\) gives
\[
\sum_{A\in\bigcup_{\ell\in L}C_\ell}p^{|A|}q^{|U|-|A|}
\le
p^{u_R}
\sum_{\substack{S\subseteq U_B\\ B\subseteq S}}
  p^{|S|}q^{|U_B|-|S|}.
\]
The set \(B\) is contained in \(U_B\) and has size \(u_B\).  Summing over the
free blue coordinates in \(U_B\setminus B\) gives
\[
\sum_{\substack{S\subseteq U_B\\ B\subseteq S}}
  p^{|S|}q^{|U_B|-|S|}
=p^{|B|}=p^{u_B}.
\]
Combining the last two displays yields the desired bound
\[
\sum_{A\in\bigcup_{\ell\in L}C_\ell}p^{|A|}(1-p)^{|U|-|A|}
\le p^{u_R}p^{u_B}.
\]
\end{proof}
